\usepackage[top=3.5cm,right=3cm,bottom=4cm,left=3cm]{geometry}  
\usepackage{amsmath,amsfonts,amssymb,amsthm}
 \usepackage[marginal,stable,bottom]{footmisc}     % for footnotes: marginal --> the same margins as text, 
                                                                                  %                       stable--> ?
                                                                                  %                       bottom --> starting the footnotes at a fixed place at the bottom of the page.
 
%\usepackage{perpage}                                             % for footnotes: starting from 1 perpage
%\usepackage{zref-abspage}
%\zmakeperpage{footnote}
%\MakePerPage{footnote}
%added by samim zarrin pour : DO NOT UNCOMMENT ABOVE LINES (4)
\usepackage{zref-perpage}
\zmakeperpage{footnote} 

\usepackage{cite}                                                     % for collecting citations: [1,2,3,4] --> [1-4]
\usepackage{setspace}                                            % for switching between double/single space in document
\allowdisplaybreaks                                                    % breaking the lines&pages when needed for the style.
\usepackage{parskip}                                               % ? 

%\relpenalty=9999                                            % show the neccessity of breaks for lines, changing the number to 10000 cause no break in lines.
%\binoppenalty=9999

\usepackage{xecolour}
\usepackage{makeidx}
\usepackage{verbatim}
\usepackage[colorlinks,linkcolor=blue,citecolor=blue]{hyperref}
\usepackage{graphicx}
\usepackage{ifthen}

 \parindent0pt


\makeatletter
\pdfstringdefDisableCommands{%
\let\lr\@firstofone
}
\makeatother
% ----------------------------------------------------------------------------
%added by samim zarrin pour
%to include hyperlinks
\usepackage{hyperref}
%to highlight freefem++ code
\usepackage{ff++listings}
%for figures
\usepackage{caption}
\usepackage{subcaption}
%for R symbol
\usepackage{amssymb}
%for algorithms
\usepackage[]{algorithm2e}

\usepackage[extrafootnotefeatures]{xepersian}

\usepackage[Lenny]{fncychap}

\settextfont[Scale=1.07]{XB Niloofar}
\setlatintextfont[Scale=1.05]{Times New Roman}
\setdigitfont[Scale=1.05]{Yas} 
% قلم برای اعداد به صورت فارسی با صفر توخالی، در صورتی که بخواهیم اعداد انگلیسی نوشته شوند این خط را غیر فعال می‌کنیم.

\defpersianfont\nastaliq[Scale=2.0]{IranNastaliq}
% قلم برای نوشتن تقدیم
\defpersianfont\nastaliqone[Scale=1.0]{IranNastaliq}
% قلم نستعلیق با سایز متناسب با متن در صورت نیاز
\defpersianfont\anotherfont[Scale=1.2]{XP Ziba}
% قلم برای نوشتن تشکر (قلم فانتزی)
\newenvironment{fantezi}
{\anotherfont }


\makeindex


\def\beginto{
\newpage
\begin{RTL}
\begin{Huge}
\nastaliq

\begin{center}
\vspace*{0.15cm}
تقدیم به 
}

\def\endto{
~
\end{center}
\end{Huge}
\end{RTL}
}

\def\beginthanks{
\newpage

{\centering\Huge{\nastaliqone{
تشکر وقدردانی ~\\
~\\}}}
}

\def\endthanks{
~
}

\def\thanks{
\beginthanks
پس از تشکر از فداکاری‌های خانواده عزیزم و حمایت‌های مستمرشان لازم می‌دانم از دکتر هاشم صابر نجفی عزیز  به خاطر هدایت بنده به سمت مدل‌سازی و ریاضی کاربردی و استاد راهنمای صبورم، دکتر خدیجه ندائی اصل به خاطر حمایت‌های علمی و درک بسیار بالای ایشان تشکر کنم. و از هرچه بگذربم سخن دوست خوش‌تر است ...



\endthanks
}

\def\startpage{
\newpage
\vspace*{3cm}
\begin{center}
\includegraphics[width=12cm]{besm}
\end{center}
}
\makeatletter
\def\@makechapterhead#1{%
  \vspace*{50\p@}%
  {\parindent \z@ \centering\normalfont
    \ifnum \c@secnumdepth >\m@ne
      \if@mainmatter
        \huge\bfseries \@chapapp\space \tartibi{chapter} 
        \par\nobreak
        \vskip 20\p@
      \fi
    \fi
    \interlinepenalty\@M
    \Huge \bfseries #1\par\nobreak
    \vskip 40\p@
  }}
\def\@makeschapterhead#1{%
  \vspace*{50\p@}%
  {\parindent \z@ \centering
    \normalfont
    \interlinepenalty\@M
    \Huge \bfseries  #1\par\nobreak
    \vskip 40\p@
  }}
\makeatother



\renewcommand\bibname{مراجع}
\def\contentsname{فهرست}

% some extra diffinitions %%%%%%%%%%%%%%%%%
%%%%%%%%%%%%%%%%%%%%%%%%%%%%%

\def\bea{\begin{eqnarray}}
\def\eea{\end{eqnarray}}

\def\ba{\begin{array}}
\def\ea{\end{array}}

\def\ni{\noindent}
\def\nn{\nonumber}

\def\bc{\begin{center}}
\def\ec{\end{center}}


